Training data contamination inflates Large Language Model (LLM) benchmark scores, yet existing detection methods flag items without diagnosing domain-level effects or producing calibrated corrections. Systematic evaluation of twenty mitigation strategies shows none restores clean performance once contamination has occurred, motivating post-hoc diagnostic tools. We introduce ContamCal, a two-stage framework that fuses detection signals via training-free rank fusion into continuous per-item exposure scores, then embeds these scores in a contamination-augmented Item Response Theory (IRT) model with per-domain contamination parameters. Applied to MMLU with controlled contamination injection across 24 models, ContamCal reveals that contamination effects are strongly heterogeneous across cognitive domains: Math ($\gamma = +0.97$) and Science ($\gamma = +0.84$) show significant inflation while opposite-sign non-STEM effects produce a non-significant pooled estimate through dilution masking. We further document that Min-K\%++ exhibits directional reversal at small model scales (Cohen's $d = -1.55$ to $-1.98$), producing actively misleading conclusions rather than merely failing. Simulation studies establish that reliable parameter recovery requires signal AUC $\geq 0.55$ and ensemble size $J \geq 50$, well above current evaluation practice. IRT ability estimation provides calibration; domain-level $\gamma$ parameters provide diagnosis that aggregate contamination analyses cannot deliver.
